\documentclass[12pt]{article}

% الهوامش
\usepackage{geometry}
\geometry{margin=0.8in}

% تنسيقات عامة
\usepackage{xcolor}
\usepackage{fancyhdr}
\usepackage{setspace}
\usepackage[inline]{enumitem}    % enumerate* للسطر الواحد

% دعم لغات/اتجاهات
\usepackage{fontspec}
\usepackage{polyglossia}
\usepackage{bidi}                % LTR/RTL

% رياضيات
\usepackage{amsmath, amssymb, mathtools}

% ضبط اللغات والخطوط
\setmainlanguage[numerals=western]{arabic}   % افتراضي: عربي (أرقام 0-9)
\setotherlanguage{english}
\newfontfamily\arabicfont[Script=Arabic,Scale=1.05]{Amiri}
\newfontfamily\englishfont{Times New Roman}

% ترويسة/تذييل
\pagestyle{fancy}
\fancyhf{}
\renewcommand{\headrulewidth}{0pt}
\cfoot{\textenglish{Page } \thepage}

% ===== ماكرو: شارة "Question n" (يسار-إلى-يمين دائماً) =====
\newcommand{\qtitle}[1]{%
  \par\noindent\begin{LTR}\textenglish{\textbf{Question #1}}\end{LTR}\par\vspace{0.25em}%
}

% ===== ماكرو: عنصر خيار باتجاه مناسب تلقائياً =====
% فكرة: إن كان النص يحوي أحرفًا عربية نعرضه داخل \textarabic{...} حتى في سياق LTR،
% ليُحافظ على ترتيب الكلمات والرموز (Insertion/Quick داخل أقواس مثلاً).
\makeatletter
\newcommand{\AutoItem}[1]{%
  % نفترض أن PHP قد لفّ النصوص العربية بـ \textarabic{...} مسبقاً قدر الإمكان.
  % ولكن لو وصل نص غير ملفوف وفيه عربي، نعالجه هنا بخطّة بسيطة:
  \edef\@tmp{#1}%
  \ifnum0%
    \ifnum\pdfstrcmp{\detokenize{#1}}{\detokenize{\textarabic{#1}}}=0 % لو هو أصلاً ملفوف
      1\fi
    \ifnum\pdfstrcmp{\detokenize{#1}}{\detokenize{#1}}=0 % dummy لتفادي فراغ if
      0\fi
  >0 %
    #1%
  \else
    % نتركه كما هو؛ الاتجاه الخارجي سيكون LTR في قائمة الخيارات
    #1%
  \fi
}
\makeatother

% ===== ماكرو: صفّ الخيارات أفقيًا مع لف تلقائي =====
\newcommand{\renderchoices}[5]{%
  % القائمة نفسها تُجبر على LTR ليبقى A./B./... وأرقام/نقاط إنكليزية سليمة
  \begin{LTR}\begin{enumerate*}[
    label=\textbf{\Alph*.},      % A. B. C. ...
    itemjoin=\hspace{2em},
    itemsep=0.25em,
    leftmargin=0pt
  ]
    \ifx\relax#1\relax\else\item \AutoItem{#1}\fi
    \ifx\relax#2\relax\else\item \AutoItem{#2}\fi
    \ifx\relax#3\relax\else\item \AutoItem{#3}\fi
    \ifx\relax#4\relax\else\item \AutoItem{#4}\fi
    \ifx\relax#5\relax\else\item \AutoItem{#5}\fi
  \end{enumerate*}\end{LTR}\par
}

% ======= المستند =======
\begin{document}

% ===== Header (مطابق للصورة تقريباً) =====
\noindent
\begin{tabular*}{\textwidth}{@{\extracolsep{\fill}} l c r}
\textenglish{\textbf{Date:} {\begin{LTR}{\textenglish{September 5, 2025}}\end{LTR}}} &
{\Large \textbf{{\textenglish{Algorithms -1-}}}} &
\textenglish{University Of Aleppo} \\
\textenglish{\textbf{Student Name:}} &
\textenglish{\textbf{Model:} {\begin{LTR}{\textenglish{A}}\end{LTR}}} &
\textenglish{\textbf{Duration:}\ {\begin{LTR}{\textenglish{90 minuts}}\end{LTR}}} \\
\textenglish{\textbf{Student Number:}} & &
\end{tabular*}

\vspace{1.0cm}

% عنوان قسم الأسئلة LTR
{\begin{LTR}\textenglish{\textbf{Select the correct answer for each question:}}\end{LTR}}

\setstretch{1.25}

\noindent \qheader{1} {\textarabic{احسب مشتق التابع }}\( f(x) = x^2\){\textarabic{ وأوجد نهايته عند اللانهاية }}\(\lim_{x \to \infty}f(x)\)

\renderchoices{\begin{LTR}{\textenglish{14}}\end{LTR}}{\begin{LTR}{\textenglish{15}}\end{LTR}}{\underline{\begin{LTR}{\textenglish{16}}\end{LTR}}}{\begin{LTR}{\textenglish{17}}\end{LTR}}{\begin{LTR}{\textenglish{17}}\end{LTR}}
{\begin{LTR}\hfill\textenglish{[1 Point]}\end{LTR}}\par\vspace{0.5em}

\noindent \qheader{2} {\textarabic{ما هي الخوارزمية التي تستخدمها لتنظيم جوربك الملونة حسب اللون، ولكن فقط إذا كان لديك 3 جوارب؟}}

\renderchoices{{\textarabic{خوارزمية الفرز السريع (Quick Sort)}}}{{\textarabic{ خوارزمية الفرز بالإدراج (Insertion Sort)}}}{\underline{\begin{LTR}{\textenglish{0}}\end{LTR}}}{\begin{LTR}{\textenglish{0}}\end{LTR}}{\begin{LTR}{\textenglish{0}}\end{LTR}}
{\begin{LTR}\hfill\textenglish{[5 Points]}\end{LTR}}\par\vspace{0.5em}

\noindent \qheader{3} \begin{LTR}{\textenglish{Calculate the following Equations:
}}\[\begin{pmatrix} x & y & z \\ 2x & y & -z \\ x & -y & 2z  \end{pmatrix}\]{\textenglish{ = }}\[\begin{pmatrix} 6  \\ 1  \\ 5 \end{pmatrix}\]\end{LTR}

\renderchoices{\begin{LTR}{\textenglish{5}}\end{LTR}}{\underline{\begin{LTR}{\textenglish{25}}\end{LTR}}}{\begin{LTR}{\textenglish{-5}}\end{LTR}}{\begin{LTR}{\textenglish{-25}}\end{LTR}}{\begin{LTR}{\textenglish{0}}\end{LTR}}
{\begin{LTR}\hfill\textenglish{[5 Points]}\end{LTR}}\par\vspace{0.5em}

\noindent \qheader{4} \begin{LTR}{\textenglish{Find the inverse of the following matrix: }}\(\begin{bmatrix} 1 & 3 & 4 \\ 5 & 2 & 1 \\ 6 & 8 & 1 \end{bmatrix}\)\end{LTR}

\renderchoices{\underline{\begin{LTR}{\textenglish{There is no inverse for this matrix}}\end{LTR}}}{\begin{LTR}\(\begin{pmatrix} 1 & 2 & 3 \\ 4 & 5 & 6 \\ 7 & 8 & 9  \end{pmatrix}\)\end{LTR}}{\begin{LTR}{\textenglish{0}}\end{LTR}}{\begin{LTR}{\textenglish{0}}\end{LTR}}{}
{\begin{LTR}\hfill\textenglish{[4 Points]}\end{LTR}}\par\vspace{0.5em}

\noindent \qheader{5} \begin{LTR}{\textenglish{Solve for }}\(x \){\textenglish{ in the equation }}\(2x+6 = 0 \)\end{LTR}

\renderchoices{\begin{LTR}{\textenglish{-1}}\end{LTR}}{\begin{LTR}{\textenglish{-2}}\end{LTR}}{\underline{\begin{LTR}{\textenglish{-3}}\end{LTR}}}{\begin{LTR}{\textenglish{3}}\end{LTR}}{}
{\begin{LTR}\hfill\textenglish{[1 Point]}\end{LTR}}\par\vspace{0.5em}

\noindent \qheader{6} \begin{LTR}\(x=3+6 \){\textenglish{ what is the value of x?}}\end{LTR}

\renderchoices{\begin{LTR}{\textenglish{7}}\end{LTR}}{\underline{\begin{LTR}{\textenglish{9}}\end{LTR}}}{\begin{LTR}{\textenglish{8}}\end{LTR}}{\begin{LTR}{\textenglish{10}}\end{LTR}}{}
{\begin{LTR}\hfill\textenglish{[1 Point]}\end{LTR}}\par\vspace{0.5em}

\noindent \qheader{7} {\textarabic{إذا كان لديك قائمة طويلة من أسماء الفواكه، لكنك تبحث عن "تفاحة". أي خوارزمية بحث ستستخدمها إذا كنت كسولاً وتتوقع أن "تفاحة" ستكون في مكان ما في منتصف القائمة؟}}

\renderchoices{{\textarabic{ البحث في الاتساع أولاً (Breadth-First Search)}}}{{\textarabic{ البحث الثنائي (Binary Search)}}}{\underline{\begin{LTR}{\textenglish{0}}\end{LTR}}}{\begin{LTR}{\textenglish{0}}\end{LTR}}{\begin{LTR}{\textenglish{0}}\end{LTR}}
{\begin{LTR}\hfill\textenglish{[5 Points]}\end{LTR}}\par\vspace{0.5em}

\noindent \qheader{8} {\textarabic{إذا كانت لدينا خوارزمية تستغرق نفس الوقت لحل مشكلة بغض النظر عن حجمها، فما هو اسمها؟}}

\renderchoices{\begin{LTR}{\textenglish{O(n)}}\end{LTR}}{\begin{LTR}{\textenglish{O(}}\(n^{2}\){\textenglish{)}}\end{LTR}}{\underline{\begin{LTR}{\textenglish{0}}\end{LTR}}}{\begin{LTR}{\textenglish{0}}\end{LTR}}{}
{\begin{LTR}\hfill\textenglish{[1 Point]}\end{LTR}}\par\vspace{0.5em}

\noindent \qheader{9} {\textarabic{أنت في مطعم وطلبت طبقاً. بعد أن تأكل الطبق، يأتي النادل ويضع طبقًا آخر فوقه. ما هي بنية البيانات التي تشبه هذا الموقف؟}}

\renderchoices{{\textarabic{مصفوفة (Array)}}}{\underline{{\textarabic{ مكدس (Stack)}}}}{\begin{LTR}{\textenglish{0}}\end{LTR}}{\begin{LTR}{\textenglish{0}}\end{LTR}}{\begin{LTR}{\textenglish{0}}\end{LTR}}
{\begin{LTR}\hfill\textenglish{[1 Point]}\end{LTR}}\par\vspace{0.5em}

\noindent \qheader{10} \begin{LTR}{\textenglish{Solve for x..: }}\[2^x + 16 = 0 \]\end{LTR}

\renderchoices{\begin{LTR}{\textenglish{4}}\end{LTR}}{\begin{LTR}{\textenglish{-4}}\end{LTR}}{\begin{LTR}{\textenglish{2}}\end{LTR}}{\begin{LTR}{\textenglish{-2}}\end{LTR}}{\underline{\begin{LTR}{\textenglish{0}}\end{LTR}}}
{\begin{LTR}\hfill\textenglish{[5 Points]}\end{LTR}}\par\vspace{0.5em}



\end{document}
